% !TeX program = lualatex
\documentclass[11pt,a4paper]{article}
\usepackage[ngerman]{babel}
\usepackage{fontspec}
\usepackage[top=1.5cm,bottom=1.5cm,left=2cm,right=2cm]{geometry}
\usepackage{tcolorbox}
\tcbuselibrary{skins,breakable}
\usepackage{enumitem}
\usepackage{array}
\usepackage{booktabs}
\usepackage{multirow}
\usepackage{tikz}
\usetikzlibrary{arrows.meta,positioning,decorations.pathmorphing,decorations.pathreplacing}
\usepackage{siunitx}
\sisetup{locale=DE,per-mode=symbol}
\usepackage{graphicx}
\usepackage{lastpage}

% Farben
\definecolor{physikblau}{RGB}{44,90,160}
\definecolor{waermerot}{RGB}{211,47,47}
\definecolor{merkgruen}{RGB}{42,122,75}
\definecolor{hellgrau}{RGB}{245,245,245}
\definecolor{loesung}{HTML}{c62828}

% Lösungs-Befehl
\newcommand{\lsg}[1]{\textcolor{loesung}{\textbf{#1}}}

% Box-Definitionen (druckfreundlich)
\newtcolorbox{aufgabenbox}[1][]{
    colback=white,
    colframe=waermerot,
    coltitle=black,
    colbacktitle=white,
    fonttitle=\bfseries,
    boxrule=1pt,
    arc=0pt,
    left=8pt, right=8pt, top=6pt, bottom=6pt,
    title={#1}
}

\newtcolorbox{merkbox}[1][Merke]{
    colback=white,
    colframe=merkgruen,
    coltitle=black,
    colbacktitle=white,
    fonttitle=\bfseries\small,
    boxrule=1pt,
    arc=0pt,
    left=8pt, right=8pt, top=6pt, bottom=6pt,
    title={#1}
}

\newtcolorbox{protokollbox}[1][]{
    colback=white,
    colframe=physikblau,
    boxrule=0.5pt,
    arc=0pt,
    left=5pt, right=5pt, top=5pt, bottom=5pt,
    #1
}

% Lücken (für Darstellung im ML nicht verwendet, aber definiert)
\newcommand{\luecke}[1][3cm]{\underline{\hspace{#1}}}
\newcommand{\lueckelang}[1][6cm]{\underline{\hspace{#1}}}

% Niveau-Sterne
\newcommand{\niveauI}{$\star$}
\newcommand{\niveauII}{$\star\star$}
\newcommand{\niveauIII}{$\star\star\star$}

% Kopfzeile
\usepackage{fancyhdr}
\pagestyle{fancy}
\fancyhf{}
\fancyhead[L]{\small Physik Klasse 8}
\fancyhead[C]{\small Wärmelehre: Grundlagen}
\fancyhead[R]{\small\textcolor{loesung}{\textbf{MUSTERLÖSUNG}}}
\fancyfoot[C]{\small Seite \thepage\ von \pageref{LastPage}}
\renewcommand{\headrulewidth}{0.4pt}

\setlength{\parindent}{0pt}
\setlength{\parskip}{6pt}

\begin{document}

% Titelbereich
\begin{center}
    {\Large\bfseries Arbeitsblatt: Wärmelehre -- Grundlagen}\\[3pt]
    {\large\textcolor{loesung}{\textbf{--- MUSTERLÖSUNG ---}}}\\[5pt]
    {\normalsize Teilchenmodell $\bullet$ Aggregatzustände $\bullet$ Temperatur $\bullet$ Phasenübergänge}
\end{center}

\vspace{5pt}

% ===== AUFGABE 1 =====
\begin{aufgabenbox}[Aufgabe 1: Einstieg -- Wärme im Alltag \hfill \niveauI]
Notiere drei Situationen aus deinem Alltag, in denen etwas warm oder kalt wird.

\begin{enumerate}[label=\alph*)]
    \item \lsg{Heiße Herdplatte erwärmt den Topf}
    \item \lsg{Eiswürfel schmilzt im Getränk}
    \item \lsg{Sonne wärmt die Haut}
\end{enumerate}

Was haben alle diese Situationen gemeinsam?

\lsg{Es wird Energie (Wärme) von einem wärmeren zu einem kälteren Körper übertragen.}
\end{aufgabenbox}

% ===== AUFGABE 2 =====
\begin{aufgabenbox}[Aufgabe 2: Das Teilchenmodell \hfill \niveauI]
Vervollständige die drei Grundaussagen des Teilchenmodells:

\begin{enumerate}[label=\arabic*.]
    \item Alle Stoffe bestehen aus \lsg{kleinsten Teilchen}.
    \item Diese Teilchen sind in \lsg{ständiger} Bewegung.
    \item Zwischen den Teilchen gibt es \lsg{Anziehungskräfte}.
\end{enumerate}
\end{aufgabenbox}

\begin{merkbox}[Merke: Teilchenmodell]
Die Bewegung der Teilchen nennt man \textbf{thermische Bewegung}. Sie hört nie auf -- selbst wenn ein Gegenstand völlig ruhig erscheint, bewegen sich die Teilchen darin ständig!
\end{merkbox}

% ===== AUFGABE 3 =====
\begin{aufgabenbox}[Aufgabe 3: Die drei Aggregatzustände \hfill \niveauI/\niveauII]
Zeichne die Teilchenanordnung für jeden Aggregatzustand und beschreibe die Bewegung.

\begin{center}
\begin{tabular}{|c|c|c|}
\hline
\textbf{FEST} & \textbf{FLÜSSIG} & \textbf{GASFÖRMIG} \\
\hline
% === FEST: Regelmäßiges Gitter, Kugeln berührend ===
\begin{tikzpicture}
\draw[thick] (0,0) rectangle (3,3);
\foreach \x in {0.5,1.0,1.5,2.0,2.5} {
    \foreach \y in {0.5,1.0,1.5,2.0,2.5} {
        \shade[ball color=physikblau!60] (\x,\y) circle (0.24);
    }
}
\end{tikzpicture}
&
% === FLÜSSIG: Dicht/berührend aber ungeordnet ===
\begin{tikzpicture}
\draw[thick] (0,0) rectangle (3,3);
% Ungeordnet, eng beieinander, teilweise berührend
\shade[ball color=physikblau!60] (0.55,0.40) circle (0.24);
\shade[ball color=physikblau!60] (1.05,0.50) circle (0.24);
\shade[ball color=physikblau!60] (1.60,0.35) circle (0.24);
\shade[ball color=physikblau!60] (2.15,0.55) circle (0.24);
\shade[ball color=physikblau!60] (2.55,0.35) circle (0.24);
\shade[ball color=physikblau!60] (0.40,0.90) circle (0.24);
\shade[ball color=physikblau!60] (0.95,1.00) circle (0.24);
\shade[ball color=physikblau!60] (1.50,0.85) circle (0.24);
\shade[ball color=physikblau!60] (2.00,1.05) circle (0.24);
\shade[ball color=physikblau!60] (2.50,0.90) circle (0.24);
\shade[ball color=physikblau!60] (0.60,1.45) circle (0.24);
\shade[ball color=physikblau!60] (1.15,1.50) circle (0.24);
\shade[ball color=physikblau!60] (1.75,1.40) circle (0.24);
\shade[ball color=physikblau!60] (2.30,1.50) circle (0.24);
\shade[ball color=physikblau!60] (0.45,1.95) circle (0.24);
\shade[ball color=physikblau!60] (1.00,2.05) circle (0.24);
\shade[ball color=physikblau!60] (1.55,1.90) circle (0.24);
\shade[ball color=physikblau!60] (2.10,2.00) circle (0.24);
\shade[ball color=physikblau!60] (2.55,1.95) circle (0.24);
% Kurze Bewegungspfeile (gleiten umeinander)
\draw[-{Stealth[length=2mm]}, gray, semithick] (0.55,0.40) -- (0.75,0.60);
\draw[-{Stealth[length=2mm]}, gray, semithick] (1.50,0.85) -- (1.30,1.05);
\draw[-{Stealth[length=2mm]}, gray, semithick] (2.30,1.50) -- (2.45,1.30);
\draw[-{Stealth[length=2mm]}, gray, semithick] (1.00,2.05) -- (1.20,1.85);
\draw[-{Stealth[length=2mm]}, gray, semithick] (2.10,2.00) -- (1.90,2.15);
\end{tikzpicture}
&
% === GAS: Wenige Kreise, frei beweglich ===
\begin{tikzpicture}
\draw[thick] (0,0) rectangle (3,3);
% Wenige Teilchen, weit verteilt
\shade[ball color=physikblau!60] (0.5,0.5) circle (0.24);
\shade[ball color=physikblau!60] (2.4,0.8) circle (0.24);
\shade[ball color=physikblau!60] (1.3,1.6) circle (0.24);
\shade[ball color=physikblau!60] (0.6,2.5) circle (0.24);
\shade[ball color=physikblau!60] (2.2,2.3) circle (0.24);
% Lange Bewegungspfeile (schnelle Bewegung)
\draw[-{Stealth[length=2.5mm]}, gray, thick] (0.5,0.5) -- (1.1,1.1);
\draw[-{Stealth[length=2.5mm]}, gray, thick] (2.4,0.8) -- (1.8,0.3);
\draw[-{Stealth[length=2.5mm]}, gray, thick] (1.3,1.6) -- (1.9,2.2);
\draw[-{Stealth[length=2.5mm]}, gray, thick] (0.6,2.5) -- (1.2,2.1);
\draw[-{Stealth[length=2.5mm]}, gray, thick] (2.2,2.3) -- (1.6,2.7);
\end{tikzpicture}
\\
\hline
\begin{minipage}{3cm}\centering\small
Teilchen\\[3pt]
\lsg{schwingen}\\[3pt]
um feste Plätze
\end{minipage}
&
\begin{minipage}{3cm}\centering\small
Teilchen\\[3pt]
\lsg{gleiten}\\[3pt]
aneinander vorbei
\end{minipage}
&
\begin{minipage}{3cm}\centering\small
Teilchen\\[3pt]
\lsg{fliegen schnell}\\[3pt]
frei umher
\end{minipage}
\\
\hline
\end{tabular}
\end{center}

\textbf{Bindung:} Je \lsg{stärker} die Bindungskräfte zwischen den Teilchen, desto \lsg{dichter} sind die Teilchen beieinander.
\end{aufgabenbox}

% ===== AUFGABE 4 =====
\begin{aufgabenbox}[Aufgabe 4: Was ist Temperatur? \hfill \niveauII]
Erkläre in eigenen Worten, was Temperatur auf Teilchenebene bedeutet.

\begin{protokollbox}
\lsg{Die Temperatur gibt an, wie schnell sich die Teilchen eines Stoffes im Durchschnitt bewegen. Je höher die Temperatur, desto schneller bewegen sich die Teilchen und desto größer ist ihre kinetische Energie.}
\end{protokollbox}

Vervollständige:
\begin{itemize}[leftmargin=*]
    \item Hohe Temperatur $\rightarrow$ Teilchen bewegen sich \lsg{schnell} $\rightarrow$ \lsg{hohe} kinetische Energie
    \item Niedrige Temperatur $\rightarrow$ Teilchen bewegen sich \lsg{langsam} $\rightarrow$ \lsg{geringe} kinetische Energie
\end{itemize}
\end{aufgabenbox}

\begin{merkbox}[Merke: Temperatur]
\textbf{Temperatur} ist ein Maß für die \lsg{mittlere kinetische Energie} der Teilchen eines Stoffes.
\end{merkbox}

% ===== AUFGABE 5 =====
\begin{aufgabenbox}[Aufgabe 5: Temperaturskalen -- Celsius und Kelvin \hfill \niveauI/\niveauII]

\textbf{a) Fixpunkte der Celsius-Skala:}
\begin{itemize}[leftmargin=*]
    \item \SI{0}{\degreeCelsius} = \lsg{Gefrierpunkt / Schmelzpunkt von Wasser}
    \item \SI{100}{\degreeCelsius} = \lsg{Siedepunkt von Wasser (bei Normaldruck)}
\end{itemize}

\textbf{b) Die Kelvin-Skala:}

Der \textbf{absolute Nullpunkt} liegt bei \lsg{0} K = \lsg{$-$273} °C.

Bei dieser Temperatur bewegen sich die Teilchen (fast) \lsg{gar nicht mehr}.

\textbf{c) Umrechnung:}

\begin{center}
\fbox{\Large $T$ / K = $\vartheta$ / °C + 273}
\end{center}

\textbf{d) Rechne um:}
\begin{enumerate}[label=\roman*)]
    \item \SI{25}{\degreeCelsius} = \lsg{298} K
    \item \SI{350}{K} = \lsg{77} °C
    \item \SI{-40}{\degreeCelsius} = \lsg{233} K
    \item \SI{0}{K} = \lsg{$-$273} °C
\end{enumerate}
\end{aufgabenbox}

% ===== AUFGABE 6 =====
\begin{aufgabenbox}[Aufgabe 6: Erweiterung -- Fahrenheit \hfill \niveauIII]
In den USA wird die Fahrenheit-Skala verwendet.

\smallskip
\begin{center}
\fbox{\;\; $\displaystyle \vartheta_F = \frac{9}{5} \cdot \vartheta_C + 32$ \qquad\qquad $\displaystyle \vartheta_C = \frac{5}{9} \cdot (\vartheta_F - 32)$ \;\;}
\end{center}
\smallskip

\textbf{a)} Wasser gefriert bei \lsg{32} °F und siedet bei \lsg{212} °F.

\textbf{b)} Bei welcher Temperatur sind Celsius und Fahrenheit gleich? \lsg{$-$40}

\textbf{c)} Recherche: Warum hat sich Fahrenheit diese ungewöhnlichen Fixpunkte ausgedacht?

\begin{protokollbox}
\lsg{Fahrenheit wählte als 0°F die tiefste Temperatur, die er mit einer Eis-Salz-Mischung erzeugen konnte, und als 96°F die Körpertemperatur eines gesunden Menschen. So vermied er negative Werte im Alltag.}
\end{protokollbox}
\end{aufgabenbox}

% ===== AUFGABE 7 =====
\begin{aufgabenbox}[Aufgabe 7: Die sechs Phasenübergänge \hfill \niveauI/\niveauII]
Beschrifte die Pfeile mit den richtigen Begriffen:

\begin{center}
\begin{tikzpicture}[
    box/.style={draw, thick, minimum width=2.2cm, minimum height=1cm, align=center, rounded corners=4pt},
    arr/.style={-{Stealth[length=3mm]}, thick}
]
% Boxes (triangular: gas top, fest bottom-left, flüssig bottom-right)
\node[box, fill=orange!10] (gas) at (3.5,4) {GASFÖRMIG};
\node[box, fill=blue!10] (fest) at (0,0) {FEST};
\node[box, fill=green!10] (fluessig) at (7,0) {FLÜSSIG};

% Bottom: Schmelzen / Erstarren
\draw[arr, waermerot] ([yshift=3pt]fest.east) -- node[above, font=\small] {\lsg{Schmelzen}} ([yshift=3pt]fluessig.west);
\draw[arr, physikblau] ([yshift=-3pt]fluessig.west) -- node[below, font=\small] {\lsg{Erstarren}} ([yshift=-3pt]fest.east);

% Left: Sublimieren / Resublimieren
\draw[arr, purple, bend left=8] (fest.north) to node[left, font=\small] {\lsg{Sublimieren}} (gas.south west);
\draw[arr, teal, bend left=8] (gas.south west) to node[right, font=\small] {\lsg{Resublimieren}} (fest.north);

% Right: Verdampfen / Kondensieren
\draw[arr, waermerot, bend right=8] (fluessig.north) to node[right, font=\small] {\lsg{Verdampfen}} (gas.south east);
\draw[arr, physikblau, bend right=8] (gas.south east) to node[left, font=\small] {\lsg{Kondensieren}} (fluessig.north);

% Legend
\node[font=\scriptsize, text=gray] at (3.5,-0.7) {rot = Energie zuführen \quad blau = Energie abgeben};
\end{tikzpicture}
\end{center}

\textbf{Begriffe:} Schmelzen, Erstarren, Verdampfen, Kondensieren, Sublimieren, Resublimieren

\textbf{Energie:} Bei welchen Übergängen wird Energie \textbf{zugeführt}? \lsg{Schmelzen, Verdampfen, Sublimieren}

Bei welchen Übergängen wird Energie \textbf{abgegeben}? \lsg{Erstarren, Kondensieren, Resublimieren}
\end{aufgabenbox}

\newpage

% ===== AUFGABE 8 =====
\begin{aufgabenbox}[Aufgabe 8: Simulationsexperiment -- Erwärmung von Eis \hfill \niveauII]

\textbf{Simulation:} Im Lernpfad (Schritt 8) siehst du ein Becherglas mit Eis auf einer Heizplatte. Ein Thermometer zeigt die Temperatur, eine Uhr läuft mit.

\textbf{Auftrag:} Lies alle \SI{30}{s} die Temperatur ab und notiere den Messwert.

\renewcommand{\arraystretch}{1.15}
\begin{center}
\begin{tabular}{|c|p{1.5cm}||c|p{1.5cm}|}
\hline
\textbf{Zeit} & \textbf{$\vartheta$ / °C} & \textbf{Zeit} & \textbf{$\vartheta$ / °C} \\
\hline
0:00 & \lsg{$-$20} & 3:00 & \lsg{67} \\
\hline
0:30 & \lsg{0} & 3:30 & \lsg{100} \\
\hline
1:00 & \lsg{0} & 4:00 & \lsg{100} \\
\hline
1:30 & \lsg{0} & 4:30 & \lsg{100} \\
\hline
2:00 & \lsg{0} & 5:00 & \lsg{120} \\
\hline
2:30 & \lsg{33} & & \\
\hline
\end{tabular}
\end{center}

Beobachtungen während der Simulation: \lsg{Die Temperatur steigt nicht gleichmäßig. Bei 0\,°C und 100\,°C bleibt sie längere Zeit konstant (Plateau), obwohl weiter geheizt wird.}

\vspace{4pt}
\textbf{Trage deine Messwerte als Punkte in das Diagramm ein und verbinde sie:}

\begin{center}
\begin{tikzpicture}
    % Light gray background
    \fill[gray!5] (0,-1) rectangle (10,6);
    % Fine grid (5mm squares)
    \draw[step=0.5, very thin, gray!25] (0,-1) grid (10,6);
    % Major grid (1cm)
    \draw[step=1, thin, gray!50] (0,-1) grid (10,6);
    % Reference lines at 0°C and 100°C
    \draw[dashed, blue!40, semithick] (0,0) -- (10,0);
    \draw[dashed, red!40, semithick] (0,5) -- (10,5);
    \node[right, font=\tiny, blue!60] at (10.1,0) {0\,°C};
    \node[right, font=\tiny, red!60] at (10.1,5) {100\,°C};
    % Y-axis
    \draw[very thick, -{Stealth[length=3mm]}] (0,-1) -- (0,6.5);
    \node[above, font=\small] at (0,6.5) {$\vartheta$ / °C};
    % X-axis
    \draw[very thick, -{Stealth[length=3mm]}] (0,-1) -- (10.5,-1);
    \node[below, font=\small] at (5,-1.6) {Zeit $t$ / min};
    % X labels (every minute)
    \foreach \x/\lab in {0/0, 2/1, 4/2, 6/3, 8/4, 10/5} {
        \draw[thick] (\x,-1) -- (\x,-1.15);
        \node[below, font=\scriptsize] at (\x,-1.15) {\lab};
    }
    % X minor ticks (every 30s)
    \foreach \x in {1,3,5,7,9} {
        \draw (\x,-1) -- (\x,-1.08);
    }
    % Y labels (every 20°C)
    \foreach \y/\temp in {-1/-20, 0/0, 1/20, 2/40, 3/60, 4/80, 5/100, 6/120} {
        \draw[thick] (0,\y) -- (-0.15,\y);
        \node[left, font=\scriptsize] at (-0.15,\y) {\temp};
    }
    %
    % ===== LÖSUNGSKURVE =====
    \draw[loesung, very thick]
        (0,-1) -- (1,0)          % 0:00 (-20°C) → 0:30 (0°C)
        -- (4,0)                  % Schmelz-Plateau 0:30–2:00
        -- (7,5)                  % Anstieg 2:00–3:30 (100°C)
        -- (9,5)                  % Siede-Plateau 3:30–4:30
        -- (10,6);                % Anstieg 4:30–5:00 (120°C)
    % Messpunkte
    \foreach \x/\y in {0/-1, 1/0, 2/0, 3/0, 4/0, 5/1.667, 6/3.333, 7/5, 8/5, 9/5, 10/6} {
        \fill[loesung] (\x,\y) circle (3pt);
    }
    % Plateau-Beschriftung
    \node[font=\small\bfseries, text=physikblau] at (2.5,0.5) {Schmelzen};
    \draw[physikblau, thick, decorate, decoration={brace, amplitude=4pt, mirror}] (1,-0.3) -- (4,-0.3);
    \node[font=\small\bfseries, text=orange!80!black] at (8,5.5) {Sieden};
    \draw[orange!80!black, thick, decorate, decoration={brace, amplitude=4pt}] (7,5.3) -- (9,5.3);
\end{tikzpicture}
\end{center}

\textbf{a)} Markiere den Bereich, in dem die Temperatur \textbf{konstant} bleibt (Plateau). Bei welcher Temperatur liegt das Plateau? \lsg{0\,°C und 100\,°C}

\textbf{b)} Beschreibe den Verlauf deiner Kurve und erkläre, warum sie so aussieht:

\begin{protokollbox}
\lsg{Die Kurve steigt zunächst steil an ($-$20 auf 0\,°C), bleibt dann bei 0\,°C konstant (Schmelz-Plateau: Eis schmilzt, Energie wird zum Aufbrechen der Bindungen verwendet, nicht zur Erwärmung). Dann steigt die Temperatur gleichmäßig bis 100\,°C. Dort bleibt sie erneut konstant (Siede-Plateau: Wasser verdampft). Danach steigt die Temperatur des Dampfes weiter.}
\end{protokollbox}

\textit{\small Tipp: Die Messpunkte in der Simulation helfen dir beim Ablesen.}
\end{aufgabenbox}

% ===== AUFGABE 9 =====
\begin{aufgabenbox}[Aufgabe 9: Verdunsten vs. Sieden \hfill \niveauII]
Erkläre den Unterschied zwischen Verdunsten und Sieden:

\begin{center}
\renewcommand{\arraystretch}{1.5}
\begin{tabular}{|p{3cm}|p{5cm}|p{5cm}|}
\hline
& \textbf{Verdunsten} & \textbf{Sieden} \\
\hline
\textbf{Wo?} & \lsg{Nur an der Oberfläche} & \lsg{In der gesamten Flüssigkeit} \\
\hline
\textbf{Temperatur?} & \lsg{Bei jeder Temperatur} & \lsg{Nur bei Siedetemperatur (100\,°C)} \\
\hline
\textbf{Geschwindigkeit?} & \lsg{Langsam} & \lsg{Schnell, heftig (Blasen)} \\
\hline
\textbf{Beispiel:} & \lsg{Pfütze trocknet in der Sonne} & \lsg{Kochendes Wasser im Topf} \\
\hline
\end{tabular}
\end{center}

\textbf{Anwendung:} Erkläre, warum nasse Haut bei Wind kühler wird als trockene Haut:

\begin{protokollbox}
\lsg{Beim Verdunsten verlassen die schnellsten (energiereichsten) Wasserteilchen die Hautoberfläche. Dadurch sinkt die mittlere kinetische Energie der verbleibenden Teilchen -- die Haut kühlt ab (Verdunstungskälte). Wind beschleunigt die Verdunstung, weil er die feuchte Luftschicht über der Haut wegträgt.}
\end{protokollbox}
\end{aufgabenbox}

\newpage

% ===== AUFGABE 10 =====
\begin{aufgabenbox}[Aufgabe 10: Zusammenfassung -- Lückentext \hfill \niveauI]
Fülle die Lücken aus:

Alle Stoffe bestehen aus \lsg{Teilchen}, die sich \lsg{ständig} bewegen. Diese Bewegung nennt man \lsg{thermische} Bewegung.

Je nachdem, wie stark die Teilchen aneinander gebunden sind, unterscheiden wir drei \lsg{Aggregatzustände}: fest, flüssig und gasförmig.

Die \lsg{Temperatur} ist ein Maß für die mittlere \lsg{kinetische Energie} der Teilchen. Je schneller sich die Teilchen bewegen, desto \lsg{höher} ist die Temperatur.

Die \lsg{Kelvin}-Skala beginnt beim absoluten Nullpunkt (\SI{0}{K} = \SI{-273}{\degreeCelsius}). Bei dieser Temperatur bewegen sich die Teilchen (fast) gar nicht mehr.

Wenn ein Stoff seinen Aggregatzustand ändert, nennt man das \lsg{Phasenübergang}. Während dieses Vorgangs bleibt die Temperatur \lsg{konstant}, obwohl weiter Energie zugeführt wird. Diese Energie wird zum \lsg{Aufbrechen} der Bindungen zwischen den Teilchen verwendet.

\textbf{Lösungswörter:} Teilchen, ständig, thermische, Aggregatzustände, Temperatur, kinetische Energie, höher, Kelvin, Phasenübergang, konstant, Aufbrechen
\end{aufgabenbox}

% ===== ZUSATZAUFGABE =====
\begin{aufgabenbox}[Zusatzaufgabe: Transferaufgabe \hfill \niveauIII]
Ein Eiswürfel wird aus dem Gefrierfach (\SI{-18}{\degreeCelsius}) genommen und in ein Glas mit Wasser (\SI{20}{\degreeCelsius}) gelegt.

\textbf{a)} Beschreibe, was auf Teilchenebene passiert, wenn der Eiswürfel schmilzt.

\begin{protokollbox}
\lsg{Die schnelleren Wasserteilchen (20\,°C) stoßen gegen die langsameren Eisteilchen ($-$18\,°C) und übertragen dabei Energie. Die Eisteilchen schwingen immer stärker, bis die Bindungen im Kristallgitter aufbrechen. Die Teilchen werden frei beweglich -- das Eis schmilzt.}
\end{protokollbox}

\textbf{b)} Warum wird das Wasser im Glas kälter?

\begin{protokollbox}
\lsg{Die Wasserteilchen geben kinetische Energie an die Eisteilchen ab. Dadurch werden sie langsamer -- die Temperatur des Wassers sinkt.}
\end{protokollbox}

\textbf{c)} Bei welcher Temperatur würde der Schmelzvorgang aufhören, wenn man unbegrenzt viel Eis hinzufügt? Begründe.

\begin{protokollbox}
\lsg{Bei 0\,°C. Solange noch Eis vorhanden ist, wird alle zugeführte Energie zum Aufbrechen der Bindungen (Schmelzen) verwendet. Die Temperatur bleibt bei 0\,°C konstant (Schmelz-Plateau).}
\end{protokollbox}
\end{aufgabenbox}

\vfill

\begin{center}
{\small\textcolor{loesung}{\textbf{MUSTERLÖSUNG -- Nicht an Schüler verteilen!}}}
\end{center}

\end{document}
